% Section 4 — Conclusion C4: Structural Unavoidability
% Derives: AI's embedding in critical societal systems is self-reinforcing
%          and structurally unavoidable.
% Depends on: Assumption A3 (Current Sectoral Penetration) + A4 (Switching Costs)
%
% SELF-CONTAINED: All needed definitions and assumptions are stated here.

\documentclass[12pt]{article}
\usepackage{lmodern}
\usepackage{xcolor}
\usepackage{amsmath, amsthm, amssymb}
\usepackage[margin=1in]{geometry}
\input{_macros}

\newtheorem{theorem}{Theorem}
\newtheorem{lemma}{Lemma}
\newtheorem{assumption}{Assumption}
\newtheorem{definition}{Definition}

\begin{document}

\section*{Section 4: Structural Unavoidability (C4)}

\subsection*{Assumptions Used}

\begin{assumption}[Current Sectoral Penetration (A3)]
  As of 2024, AI systems are operationally embedded in the following
  \emph{critical societal systems}:
  \begin{enumerate}
    \item \textbf{Energy infrastructure}: AI manages load balancing in
      national power grids (EU, US, China).
    \item \textbf{Healthcare}: FDA-cleared AI diagnostic tools are in use at
      $>$40\% of US hospitals; AlphaFold2 is used by $>$1M researchers.
    \item \textbf{Fiscal systems}: AI drives fraud detection and tax-gap
      analysis at the IRS and HMRC.
    \item \textbf{Education}: AI-powered adaptive learning platforms serve
      $>$50M K--12 students worldwide.
  \end{enumerate}
  Each of these systems is operationally critical: failure or degradation
  causes measurable societal harm.
\end{assumption}

\begin{assumption}[Super-Linear Switching Costs (A4)]
  Removing AI from a pipeline at integration depth $d$ incurs switching cost
  $S(d) = \Omega(d^2)$.  For operationally critical systems,
  deep integration ($d \gg 1$) is the norm.
\end{assumption}

\subsection*{Statement to Prove (C4)}

\begin{theorem}[Structural Unavoidability]
  Under Assumptions A3 and A4:
  \begin{enumerate}
    \item[(i)] AI is already embedded in at least four independent critical
      societal systems at high integration depth.
    \item[(ii)] Removing AI from any of these systems would incur switching
      costs that are prohibitive relative to the system's annual operating
      budget, making removal practically infeasible.
    \item[(iii)] As usage within each system grows, integration depth
      increases, further raising switching costs — creating a
      self-reinforcing lock-in dynamic.
    \item[(iv)] Therefore, AI's role in these societal systems is
      structurally unavoidable: the societal infrastructure is now
      irrevocably AI-dependent.
  \end{enumerate}
\end{theorem}

\subsection*{Proof}

\textit{(Generator: produce a rigorous argument.  For (i), cite the four
systems in A3 and argue they are critical and independent.  For (ii), use A4
to bound switching costs relative to system budgets — you may introduce
reasonable order-of-magnitude estimates for healthcare IT, power-grid, and
education platform budgets to make the comparison concrete.  For (iii), argue
the positive-feedback loop: more use $\Rightarrow$ deeper integration
$\Rightarrow$ higher switching costs $\Rightarrow$ continued use.  For (iv),
synthesize.  Use only A3 and A4.)}

\textcolor{red}{TODO: proof of C4}

\end{document}
